% !TeX program = xelatex
\documentclass[12pt]{article}
 
% ---------- Paquetes ----------
\usepackage{fontspec}
\usepackage[spanish]{babel}
\usepackage{amsmath}
\usepackage{amssymb}
\usepackage{geometry}
\usepackage{tikz}
 
\geometry{margin=2.5cm}
 
% ---------- Datos del documento ----------
\title{\textbf{Modelo de Sensores Inerciales} \\ \large Densidad Espectral de Potencia (PSD)}
\author{Michell Porchett Lopez Plata}
\date{}
 
\begin{document}
 
\maketitle
 
\section{Modelo de Sensores Inerciales}
 
El modelo matemático que describe la medición de un acelerómetro está dado por:
 
\begin{equation}
    a_m = M_a \cdot a + b_a + n_a
\end{equation}
 
Donde cada término representa lo siguiente:
 
\begin{itemize}
    \item $a_m$: Aceleración medida.
    \item $M_a$: Matriz que representa la ganancia, desalineamiento entre ejes y escala.
    \item $a$: Aceleración real que mide el acelerómetro.
    \item $b_a$: Bias o sesgo del acelerómetro.
    \item $n_a$: Ruido de la medición.
\end{itemize}
 
\subsection{Ruido de la medición}
 
El ruido de la medición $n_a$ se modela como una variable aleatoria con distribución normal de media cero y varianza $\sigma^2$:
 
\begin{equation}
    n_a \sim N(0, \sigma^2)
\end{equation}
 
A continuación se muestra la representación gráfica de la distribución normal y su Densidad Espectral de Potencia (PSD) correspondiente:
 
\begin{center}
% --- Campana de Gauss ---
\begin{tikzpicture}[scale=1]
    \node at (0,2.1) {$N(0,\sigma^2)$};
    % Ejes
    \draw[-latex] (-2.5,0) -- (2.5,0);
    \draw[-latex] (0,0) -- (0,2);
    % Curva gaussiana
    \draw[thick, purple, domain=-2.2:2.2, samples=80, smooth]
        plot (\x, {1.8*exp(-\x*\x/2)});
    % Marcas en el eje x
    \foreach \x/\lbl in {-2/-2\sigma, -1/-\sigma, 0/0, 1/\sigma, 2/2\sigma} {
        \draw (\x,0.05) -- (\x,-0.05) node[below] {\footnotesize $\lbl$};
    }
\end{tikzpicture}
\hspace{1.5cm}
% --- PSD constante ---
\begin{tikzpicture}[scale=1]
    \node at (0,2.1) {PSD};
    % Ejes
    \draw[-latex] (-2.5,0) -- (2.5,0) node[right] {$\omega$};
    \draw (-2.3,0) node[below] {\footnotesize $-\omega$};
    \draw[-latex] (0,0) -- (0,2);
    % Linea PSD constante, tipo onda seno
    \draw[thick, purple, domain=-2.3:2.3, samples=150, smooth]
        plot (\x, {1 + 0.06*sin(deg(6*\x))});
\end{tikzpicture}
\end{center}
 
La Densidad Espectral de Potencia (PSD) de este tipo de ruido es \textbf{constante} en todo el espectro de frecuencias, y sus unidades están dadas por:
 
\begin{equation}
    PSD = \text{Constante} \; \left[ \dfrac{(\text{Unidades})^2}{Hz} \right]
\end{equation}
 
\newpage
 
\section{Ejemplo: Sensor MPU-6050}
 
 
\subsection{Densidad de ruido (Noise Density)}
 
\begin{equation}
    N_a = 400 \; \dfrac{\mu g}{\sqrt{Hz}}
\end{equation}
 
Convirtiendo las unidades de $\mu g$ a $m/s^2$ (multiplicando por $9.81 \; m/s^2$):
 
\begin{equation}
    N_a = \dfrac{400 \times 10^{-6} \cdot 9.81 \; m/s^2}{\sqrt{Hz}}
\end{equation}
 
\begin{equation}
    N_a = 3.924 \times 10^{-3} \; \dfrac{m/s^2}{\sqrt{Hz}}
\end{equation}
 
\subsection{Cálculo de la PSD}
 
Sabiendo que la relación entre la densidad de ruido y la PSD es:
 
\begin{equation}
    PSD = N_a^2
\end{equation}
 
Sustituyendo el valor obtenido:
 
\begin{equation}
    PSD = 1.539 \times 10^{-5} \; \dfrac{(m/s^2)^2}{Hz}
\end{equation}
 
\subsection{Cálculo de la varianza}
 
La varianza del ruido se relaciona con la PSD y la frecuencia de muestreo $f_s$ mediante:
 
\begin{equation}
    \sigma^2 = PSD \cdot f_s \quad \text{(} f_s \text{: frecuencia de muestreo)}
\end{equation}
 
Tomando $f_s = 1 \times 10^{3} \; Hz$ como frecuencia de muestreo:
 
\begin{equation}
    \sigma^2 = 1.539 \times 10^{-5} \cdot (1 \times 10^{3})
\end{equation}
 
\begin{equation}
    \boxed{\sigma^2 = 0.01539}
\end{equation}
 
\end{document}